% Publication fragment: standard amsmath/amsthm environments.
% Bibliography: geometry_references.bib. Audit manifest: geometry_status.txt.
\section{Projective moment geometry}
\label{geo:section}

A probability distribution of pure quantum states can be studied through its
moments, through the algebraic equations satisfied by its support, or through
the singular potentials against which it is tested. These descriptions retain
different information. The first moment specifies an average density matrix;
the second moment specifies all quadratic experiments on matched preparations;
higher moments can distinguish ensembles whose lower moments coincide. On the
geometric side, a homogeneous equation of degree \(k\) becomes a positive test
on \(k\) matched copies. We use this relation between moments and equations
in both directions.

Throughout this section, \(P_z=zz^*/\|z\|^2\), and \(\sigma\) denotes the
unitary-invariant probability measure on the relevant complex projective
space. A projective \(k\)-design is a Borel probability measure satisfying
\begin{equation}
 M_k(\mu):=\int P_z^{\otimes k}\,d\mu(z)
 =\frac{I_{\operatorname{Sym}^k\mathbb C^D}}
 {\binom{D+k-1}{k}}.
 \label{geo:design}
\end{equation}
Partial trace gives all lower moment identities. The definition allows
weighted discrete designs and continuous measures. The connection between
projective second designs and informationally complete measurements is
standard \cite{geo_scott2006}; the results below exploit the additional
geometry of a design's support.

\subsection{A sharp first-moment entropy threshold in every dimension}
\label{geo:threshold}

The geometric energy is the Legendre transform of the Monge--Amp\`ere
energy. Its normalization matters. Put \(d=D-1\), \(X=\mathbb P^d\),
\(H=\mathcal O(1)\), \(L=-K_X=DH\), and \(\omega_0=D\omega_{\rm FS}\), with
\(\int_X\omega_{\rm FS}^d=1\). For a smooth strictly
\(\omega_0\)-plurisubharmonic potential \(u\), set
\begin{align}
 \mathcal E_L(u)
 &=\frac{1}{D^{d+1}}\sum_{j=0}^{d}
       \int_Xu\,\omega_u^j\wedge\omega_0^{d-j},
 &\omega_u&=\omega_0+dd^cu,\\
 E_L(\mu)
 &=\sup_u\left\{\mathcal E_L(u)-\int_Xu\,d\mu\right\}.
 \label{geo:energy}
\end{align}
Constants cancel in the second line. A measure is \emph{balanced} if
\(M_1(\mu)=I_D/D\). Balance is simultaneously a vanishing moment-map
condition and the assertion that the mean quantum state is maximally mixed.

\begin{theorem}[Sharp balanced entropy threshold]
\label{geo:sharp}
Let \(H_D=\sum_{j=1}^D j^{-1}\). The supremum of the coefficients
\(\gamma\) for which
\[
 \operatorname{Ent}(\mu\mid\sigma)\ge\gamma E_L(\mu)-C_\gamma
\]
holds uniformly over balanced finite-entropy measures is
\begin{equation}
 \Gamma_d^{\rm bal}=\frac{d}{D-H_D}.
 \label{geo:gamma}
\end{equation}
Every strict subcritical coefficient admits such a constant; every strict
supercritical coefficient fails. In particular, the first-moment constant
on \(\mathbb P^2\) is \(12/7\). The assertion concerns the supremal
coefficient and does not require boundedness at the endpoint.
\end{theorem}

The first values are
\[
 2,\qquad \frac{12}{7},\qquad \frac{36}{23},\qquad
 \frac{240}{163},\qquad \frac{100}{71}.
\]
The harmonic number has a concrete origin: for a uniform random point
\((x_0,\ldots,x_{D-1})\) of the probability simplex,
\(\mathbb E\max_i x_i=H_D/D\). The proof converts a balanced allocation of
monomials to projective-frame points into a statement about singularities,
then into an entropy inequality. This explains why the same constant
appears in a combinatorial allocation problem and in pluripotential theory.

\fig{fig04_jet_geometry}{Minimal jet allocation and the sharp balanced
coefficient. (a) Degree-24 monomials on $\mathbb P^2$ are plotted in their
normalized exponent simplex. The color is the minimal normalized order
$1-\max_i\alpha_i/24$ among the three coordinate centers. White boundaries
separate the maximizing-coordinate regions; ties can be allocated to meet
equal quotas. The continuum average is $7/18$. (b) The resulting
anticanonically normalized coefficient
$\Gamma_d^{\rm bal}=d/(d+1-H_{d+1})$ decreases toward one as complex
dimension grows. Its projective-plane value is $12/7$; the dashed line marks
one, not a separately proved microscopic integrability threshold.}{fig:jets}

\paragraph{Balanced jets and their invariant.}
Choose \(M\) with
\[
 N=\binom{M+d}{d}=Dn,\qquad
 S_M=\sum_{\substack{\alpha\in\mathbb N^D\\|\alpha|=M}}
                  (M-\max_i\alpha_i).
\]
There are arbitrarily large such levels: \(M=1+kD\) suffices. At coordinate
point \(e_i\), the order of \(z^\alpha\) is \(M-\alpha_i\). A joint jet
determinant using \(n\) rows at each of the \(D\) centers therefore has
total order at least \(S_M\). Equality is possible: assign each monomial
equally among its maximum-exponent centers, obtaining load \(N/D\) by
symmetry, and use integral max-flow to obtain the same integer quotas.

Collect the minimal-order minors into a nonzero vector \(P_0\) in
\[
 \left(\bigotimes_{i=0}^{D-1}
 \bigwedge^n\bigoplus_{\ell=0}^{M}
              \operatorname{Sym}^{\ell}(\mathbb C^d)^*\right)_{\deg S_M}.
\]
Use the invariant monomial norms \(\|x^\beta\|^2=\beta!/|\beta|!\).
After twisting by \(\prod_i\det_i^{-S_M/(Dd)}\), this vector has zero
moment map for \(\operatorname{GL}(d)^D\). A local root operator changes
one global monomial label, so its image has no overlap with the supported
assignments. Permutations of tangent coordinates kill the traceless
diagonal expectations, and permutations of centers distribute the total
degree equally. Polynomial pullback here is a right action; transposing
each matrix converts it to a left representation without changing the
determinant character.

Kempf--Ness, after clearing the rational character by a tensor power,
supplies a homogeneous semi-invariant \(F\) of some degree \(q>0\):
\begin{equation}
 F(P_0)\ne0,\qquad
 F(\rho(g)P)=\prod_i(\det g_i)^{qS_M/(Dd)}F(P).
 \label{geo:jet-semi-invariant}
\end{equation}
Let \(\mathcal U\subset X^D\) be the locus of \(D\) linearly independent
rays. Projective transformations act transitively on \(\mathcal U\).
All lower-total-order quota minors vanish identically there, and the
semi-invariant in \eqref{geo:jet-semi-invariant} is nonzero there.

\paragraph{From jets to log-canonical thresholds.}
For a basis \(s_1,\ldots,s_N\) of \(H^0(X,MH)\), form the balanced basis
ideal on \(X^D\),
\[
 \mathfrak b_s=
 \left\langle\prod_{i=0}^{D-1}\prod_{j\in B_i}s_j(z_i):
 \bigsqcup_i B_i=\{1,\ldots,N\},\quad |B_i|=n\right\rangle.
\]
The log-canonical threshold is the largest exponent for which the inverse
squared norm of ideal generators is locally integrable. On a resolution
it is computed by log discrepancy divided by ideal order; discrepancy
records the Jacobian cost of changing coordinates.

Let \(v_i\) be a monomial valuation on a resolving chart
\(F_i:y_i\mapsto x_i\), with weights \(a_{i1},\ldots,a_{id}\ge0\), and
\(J_i=\det dF_i\). Then
\[
 A_X(v_i)=\sum_j a_{ij}+v_i(J_i).
\]
The chain rule for the minimal jet vector has no surviving nonlinear
correction: every such term contains a lower-total-order minor, which
vanishes on \(\mathcal U\). Therefore
\[
 \widetilde P=\rho(dF_i)P(F_i),\qquad
 F(\widetilde P)=
 \prod_iJ_i^{qS_M/(Dd)}F(P(F_i)).
\]
The last factor is a valuation unit. Scaling jet rows by resolving
coordinates charges their differentiation costs to the same determinant
character. Some minimal minor consequently satisfies
\[
 \text{valuation of minor}+\text{differentiation cost}
 \le\frac{S_M}{Dd}\sum_i A_X(v_i).
\]
Apply its derivatives to the evaluation determinant in independent
particle variables, and only then identify the \(n\) variables within
each cluster. Identification can increase valuation. Thus the displayed
bound gives an upper bound for the relevant product valuation of the
independent-particle determinant.

Resolve the divisors of all \(s_j\) separately in each factor. The
pullback of \(\mathfrak b_s\) is a monomial ideal on the product resolution,
so its threshold is computed by valuations of this product form. Its
order is at most the determinant order. Hence
\begin{equation}
 \operatorname{lct}_{\mathcal U}(\mathfrak b_s)
 \ge\frac{Dd}{S_M}.
 \label{geo:basis-lct}
\end{equation}
When \(S_M=0\), the ideal is the unit ideal on \(\mathcal U\).
This statement concerns balanced basis ideals; arbitrary valuations on a
full particle configuration space need not be product valuations.

\paragraph{The analytic transfer.}
Fix a compact subset of \(\mathcal U\) and one degree \(M\). Strictly
below \eqref{geo:basis-lct}, inverse ideal norms are integrable uniformly
over the compact group of orthonormal section bases, by semicontinuity
of complex singularity exponents \cite{geo_demailly_kollar2001}.
For \(Q(z)=\sum_j e^{\lambda_j}|s_j(z)|^2\), expansion gives
\[
 \prod_iQ(z_i)^n\ge e^{\sum_j\lambda_j}
 \sum_B\left|\prod_i\prod_{j\in B_i}s_j(z_i)\right|^2.
\]
This yields the quantized local inequality for every
\(\beta<dMN/S_M\). No estimate uniform in growing \(M\) is used.
The normalized monomial lattice converges to the uniform simplex law:
\begin{equation}
 \frac{S_M}{MN}\longrightarrow\frac{D-H_D}{D}.
 \label{geo:monomial-limit}
\end{equation}
Indeed, integrating
\[
 \mathbb P(\max_i x_i>a)
 =\sum_{j=1}^D(-1)^{j+1}\binom Dj(1-ja)_+^{D-1}
\]
gives \(\mathbb E\max_i x_i=H_D/D\).

For the Hilbert map
\[
 H_\psi(s,t)=\int_Xh_{\rm FS}^M(s,t)e^{-M\psi}\,d\sigma,\qquad
 v=\operatorname{FS}(H_\psi),
\]
and \(E_M(v)=(MN)^{-1}\log(\det H_0/\det H_\psi)\), one has the exact
comparison
\begin{equation}
 \mathcal E_H(\psi)\le E_M(v),\qquad
 \int_Xe^{M(v-\psi)}\,d\sigma=N.
 \label{geo:exact-quantization}
\end{equation}
Here \(\mathcal E_H\) uses \(\omega_{\rm FS}\). To justify the exact
coefficient, apply the direct-image maximum principle of
\cite[Propositions 4.1--4.2]{geo_zhang2024} along the weak geodesic
from \(0\) to \(\psi\). The auxiliary bundle \(MH-K_X\) has curvature
\[
 M(\pi^*\omega_{\rm FS}+dd^c\psi)+D\pi^*\omega_{\rm FS}\ge0,
\]
and the reference Bergman density is exactly \(N\). Differentiating the
Hilbert/Bergman geodesic comparison at the reference endpoint gives the
first inequality; the second is the trace of Hilbert orthonormality.
This exact endpoint-derivative specialization is spelled out again in
Theorem~\ref{an:lower-comparison}; it is not obtained by setting the error
parameter to zero in an asymptotic quantization bound.

Given \(\lambda<dD/(D-H_D)\), choose a larger strict-subcritical \(\beta\)
and a single sufficiently large admissible \(M\) with
\(\lambda/M+\lambda/\beta\le1\). H\"older's inequality with these weights
transfers the quantized estimate using \eqref{geo:exact-quantization}.
Since \(\mathcal E_L(D\psi)=D\mathcal E_H(\psi)\), for
\(\gamma<d/(D-H_D)\) this gives
\begin{equation}
 \gamma\mathcal E_L(u)+\frac1D\sum_{i=0}^{D-1}
          \log\int_{U_i}e^{-\gamma u}\,d\sigma
 \le C_{\gamma,U},
 \label{geo:local-entropy}
\end{equation}
whenever the \(U_i\) have disjoint closures with closed product in
\(\mathcal U\).

\paragraph{Globalization by a determinantal coupling.}
For unit representatives set
\[
 \mathfrak D(z_1,\ldots,z_D)=|\det(z_1,\ldots,z_D)|,\qquad
 c_D=\frac{D^{D-1}}{(D-1)!}.
\]
If \(\nu\) is balanced, \(c_D\mathfrak D^2\nu^{\otimes D}\) has all
marginals \(\nu\). Assume first that \(\nu\ge\epsilon\sigma\).
The conditional operator and the derivative of common marginal scaling are
\[
 Tf(z)=\frac D{D-1}\left(\int f\,d\nu-
     \operatorname{tr}\left(P_z\int fP\,d\nu\right)\right),\qquad
 \mathcal L=I+(D-1)T.
\]
On traceless Hermitian matrices put
\(\mathcal B(M)=D\int\operatorname{tr}(MP)P\,d\nu\). Haar's second
moment and the variance inequality imply
\[
 0\le\mathcal B\le\left(1-\frac{D\epsilon}{D+1}\right)I.
\]
For \(m=\int g\,d\nu\), \(G=\int g(P-I/D)\,d\nu\), the inverse is
\begin{equation}
 \mathcal L^{-1}g=g-\frac{D-1}{D}m+
 D\operatorname{tr}\bigl((I-\mathcal B)^{-1}G\,P\bigr),\qquad
 \|\mathcal L^{-1}\|_{\infty\to\infty}\le\kappa:=\frac{D+3}{\epsilon}.
 \label{geo:marginal-inverse}
\end{equation}
Truncate the kernel to \(\mathfrak D\ge\delta\), where
\(\delta=(8\kappa D\sqrt{c_D})^{-1}\). Each marginal loses at most
\(\eta=c_D\delta^2\). Write the common scaling as \(a=1+h\). Its
equation is \(-r+\mathcal Lh+R(h)=0\), with \(\|r\|_\infty\le\eta\).
On \(\|h\|_\infty\le b=2\kappa\eta\),
\[
 \operatorname{Lip}(R)\le D\eta+2D(D-1)b,\qquad
 \kappa\operatorname{Lip}(R)<\frac18.
\]
Contraction gives the exact coupling
\[
 d\pi=\prod_i a(z_i)c_D\mathfrak D^2
          \mathbf1_{\{\mathfrak D\ge\delta\}}\,d\nu^{\otimes D}.
\]
A fixed finite partition of diameter at most \(\delta/(2D)\) in
projective chordal distance localizes it into finitely many frame
products, since the determinant is separately \(1\)-Lipschitz. Each good
closed product has determinant at least \(\delta/2\), hence disjoint
coordinate-cell closures. Slight enlargements give the \(U_i\) in
\eqref{geo:local-entropy}. Conditioning on labels expresses
\(\nu=\sum_jp_j\nu_j\), where \(\nu_j\) is the equal average of \(D\)
measures in one frame product, and the number \(J\) depends only on
\(D,\epsilon\). The finite-mixture entropy identity gives
\[
 \sum_jp_j\operatorname{Ent}(\nu_j\mid\sigma)
 \le\operatorname{Ent}(\nu\mid\sigma)+\log J.
\]
Entropy duality in \eqref{geo:local-entropy} and energy convexity now
prove the required inequality for \(\nu\ge\epsilon\sigma\).

To remove the Haar floor, normalize \(\int u\,d\sigma=0\) and set
\(J(u)=-\mathcal E_L(u)\). Integration by parts gives
\(I(u)\ge(1+1/d)J(u)\) and \(c\frac{d}{dc}J(cu)=I(cu)\), so
\(J(cu)\le c^{1+1/d}J(u)\). Taking \(c=\theta^d\) in the Legendre
transform proves the extended-valued inequality
\begin{equation}
 E_L(\theta\mu+(1-\theta)\sigma)\ge\theta^{d+1}E_L(\mu).
 \label{geo:dilution}
\end{equation}
Entropy convexity transfers the estimate to \(\mu\) with coefficient
\(\gamma\theta^d\). Letting \(\theta\uparrow1\) proves every
strict-subcritical case of Theorem~\ref{geo:sharp}.

\paragraph{Explicit sharpness.}
For \(t>\log2\), set
\[
 C_i(t)=\{[z]:1-|z_i|^2\le e^{-t}\},\qquad
 \mu_t=\frac{\mathbf1_{\cup_i C_i(t)}}{De^{-dt}}\,\sigma.
\]
The disjoint caps each have mass \(e^{-dt}\). Their torus and permutation
symmetries make \(\mu_t\) balanced, and its entropy is \(dt-\log D\).
For an orthonormal monomial basis and \(r_\alpha=M-\max_i\alpha_i\),
consider
\[
 \psi_t=\frac1M\log\sum_{|\alpha|=M}e^{tr_\alpha}|s_\alpha|^2.
\]
Subgeodesic convexity gives
\(\mathcal E_H(\psi_t)\ge(\log N)/M+tS_M/(MN)\). On cap \(i\),
\(r_\alpha\le M-\alpha_i\), so the multinomial theorem gives
\(\psi_t\le(\log N)/M+\log2\). Passing \(M\to\infty\) in energy duality,
\[
 E_L(\mu_t)\ge(D-H_D)t-D\log2.
\]
Every strict-supercritical coefficient fails along this family. Unitary
convolution and a small Haar admixture produce smooth positive balanced
examples with the same limiting obstruction.

\paragraph{Relation to stability theory.}
Moment-constrained entropy coercivity is the analytic threshold in the
probabilistic approach to K\"ahler--Einstein metrics with continuous
automorphisms. In the stated anticanonical normalization,
\cite[Propositions 3.6 and 3.10]{AndreassonBermanSvensson2026} identify this balanced
threshold on projective space with the analytic polystability threshold.
The theorem calculates that analytic quantity; microscopic integrability
of an interacting particle determinant requires additional information.

\subsection{Which plane cubics can carry a second design?}
\label{geo:cubic-classification}

A second design is highly isotropic, but its support need not fill
projective space. The following classification identifies exactly how
a cubic equation can coexist with that isotropy. Three structures enter
the proof essentially: a rank inequality for matrix pencils, the
moment-map description of multipartite entanglement, and the order-three
grading of the exceptional Lie algebra \(E_6\).

Let \(\omega=e^{2\pi i/3}\), and define the nine Hesse rays
\begin{equation}
 \mathcal H=\left\{
 \frac{(0,1,-\omega^j)}{\sqrt2},
 \frac{(1,0,-\omega^j)}{\sqrt2},
 \frac{(1,-\omega^j,0)}{\sqrt2}:j=0,1,2\right\},
 \qquad
 \mu_{\mathcal H}=\frac19\sum_{h\in\mathcal H}\delta_h.
 \label{geo:hesse-points}
\end{equation}
They form a symmetric informationally complete set:
\(\operatorname{tr}(P_hP_{h'})=1/4\) for \(h\ne h'\).
Their incidence geometry and linear dependencies are discussed in
\cite{geo_hesse2009,geo_dang2013}.

\begin{theorem}[Complete cubic-support classification]
\label{geo:cubic}
For a nonzero homogeneous ternary cubic \(F\), a projective second design
supported on \(\{F=0\}\) exists if and only if, after a unitary coordinate
change and multiplication by a nonzero scalar, \(F\) is either
\[
 z_0^3+z_1^3+z_2^3-\lambda z_0z_1z_2,\qquad \lambda^3\ne27,
\]
or \(z_0z_1z_2\).
In the smooth case the measure is uniquely \(\mu_{\mathcal H}\), in the
corresponding unitary coordinates. On the coordinate triangle every such
measure is
\[
 \mu=\frac13(\nu_0+\nu_1+\nu_2),
\]
where \(\nu_i\) is supported on \(z_i=0,\ |z_j|=|z_k|\), and its
first two phase Fourier coefficients vanish.
\end{theorem}

The conclusion is unitary, rather than merely projective: a general
projective transformation changes the prescribed Hermitian moments.
Nodal or cuspidal cubics, a line plus an irreducible conic, nonorthogonal
triangles, concurrent triples, and repeated components are excluded.

\begin{proof}
\emph{Step 1: a determinant with two balanced tensor factors.}
Every ternary cubic has a \(3\times3\) linear determinantal representation
\[
 F(z)=\det M(z),\qquad M(z)=\sum_{i=0}^2z_iM_i.
 \tag{a}
\]
For a smooth cubic, use the Hesse normal form
\cite[Lemma 1]{geo_hesse2009} and the Moore matrix
\[
 \mathsf M(a,z)=
 \begin{pmatrix}
 a_0z_0&a_1z_2&a_2z_1\\
 a_2z_2&a_0z_1&a_1z_0\\
 a_1z_1&a_2z_0&a_0z_2
 \end{pmatrix},
 \quad
 \det\mathsf M=S(a)T(z)-T(a)S(z),
\]
where \(S(z)=\sum_i z_i^3\), \(T(z)=z_0z_1z_2\). Choose \(a\) on the
cubic with \(T(a)\ne0\).
For a singular irreducible cubic, put a singular point at
\([0:0:1]\), so \(F=zQ_2(x,y)+Q_3(x,y)\). If \(Q_2\) has distinct
linear factors, normalize it to \(xy\) and absorb the terms divisible
by \(xy\) into \(z\), obtaining \(xyz+ax^3+dy^3\).
If \(Q_2\) is a square, normalize it to \(y^2\), absorb its multiples
into \(z\), then translate \(x\) to obtain \(x^3-y^2z\).
The two resulting determinants are
\[
 \begin{pmatrix}x&0&y\\0&y&x\\-dy&-ax&z\end{pmatrix},
 \qquad
 \begin{pmatrix}x&y&0\\0&x&z\\-y&0&x\end{pmatrix}.
\]
A reducible cubic is a line times a conic; represent the conic by a
\(2\times2\) linear determinant and add the line as a diagonal block.
This includes repeated factors.

Regard the coefficients of \(M\) as
\(\Psi\in(\mathbb C^3)^{\otimes3}\), and put
\[
 s=\|\Psi\|^2,\quad
 R=\sum_iM_iM_i^*,\quad
 C=\sum_iM_i^*M_i,\quad
 T_{ij}=\operatorname{tr}(M_iM_j^*).
\]
Up to a transpose for the second factor, these are its one-factor
reduced density matrices, each of trace \(s\).
The coefficients of \(\det M(z)\) are invariant under
\(M\mapsto AMB^{\mathsf T}\), \(A,B\in SL_3\).
Their nonvanishing prevents zero from belonging to the orbit closure.
The closure has a nonzero point of minimum Euclidean norm. Differentiating
its norm in all traceless Hermitian directions on the first two factors
gives
\[
 R=C=\frac{s}{3}I.
 \tag{b}
\]
The determinant remains exactly \(F\) at the limit. This argument does
not assume a finite invertible scaling of the original pencil.

\emph{Step 2: the design completes the balance.}
The Hermitian quartic
\[
 q_M(z)=2\operatorname{tr}\bigl((M(z)M(z)^*)^2\bigr)
       -\bigl(\operatorname{tr}(M(z)M(z)^*)\bigr)^2
\]
is nonnegative on \(\det M=0\). Indeed, a positive matrix of rank at most
two, with possibly zero eigenvalues \(b_1,b_2\), gives
\(q_M=(b_1-b_2)^2\). Haar's second-moment contractions give
\[
 12\int q_M\,d\sigma
 =2\operatorname{tr}R^2+2\operatorname{tr}C^2
      -\operatorname{tr}T^2-s^2.
 \tag{c}
\]
If a second design is supported on the cubic, then (b), (c), and its
moment identity imply
\[
 0\le12\int q_M\,d\mu
 =-\left\|T-\frac{s}{3}I\right\|_{\rm HS}^2\le0.
\]
Thus the third marginal is scalar. The measure constraint has forced a
three-qutrit tensor to be fully balanced.

\fig{fig03_tensor_balance}{The tensor used in the cubic classification.
The three matrix slices display the nonzero coefficients of a vector in
the Cartan subspace, with illustrative values $a=1$, $b=0.65$, and $c=0.35$.
After normalization, all three one-qutrit marginals are $I_3/3$, regardless
of these coefficients. The slice determinant is
$(a^3+b^3+c^3)z_0z_1z_2-abc(z_0^3+z_1^3+z_2^3)$, a member of the Hesse
pencil. The proof first obtains full marginal balance from the design
identity and only then invokes the $E_6$ Cartan representation.}{fig:tensor}

\emph{Step 3: the exceptional-Lie-algebra input.}
For the \(SL_3^3\) action, scalar one-factor marginals characterize
critical vectors. Their orbits are closed, and two critical vectors in
one orbit are locally unitarily equivalent, by Kempf--Ness
\cite[Theorem 14]{geo_sauerwein2018}. Vinberg's Cartan theorem for the
order-three \(E_6\) representation states that every closed orbit meets
\[
 \mathfrak c=\operatorname{span}\left\{
 |000\rangle+|111\rangle+|222\rangle,\,
 |012\rangle+|120\rangle+|201\rangle,\,
 |021\rangle+|210\rangle+|102\rangle\right\}.
\]
The statement includes nongeneric semisimple orbits
\cite[Section 7.1]{geo_ditrani2023}. Every vector in this space already
has scalar marginals: the three supports are disjoint Latin squares.
The tensor from (b) is consequently related to \(a u_1+b u_2+c u_3\)
by local unitaries. Its determinant is
\[
 (a^3+b^3+c^3)z_0z_1z_2
             -abc(z_0^3+z_1^3+z_2^3).
 \tag{d}
\]
The first two local unitaries change the determinant by a scalar; the
third gives a unitary change of its variables. Its nonzero singular
members are precisely orthogonal triangles. For example,
\[
 x^3+y^3+z^3-3xyz
 =(x+y+z)(x+\omega y+\omega^2z)(x+\omega^2y+\omega z),
\]
and the three coefficient vectors are orthogonal.

\emph{Step 4: rigidity on a smooth Hesse cubic.}
Define
\[
 r(z)=\sum_i|z_i|^4-\frac12\|z\|^4,\qquad
 D(z)=\sum_i z_i^2\overline{z_{i+1}z_{i+2}}.
\]
The Moore matrix satisfies
\[
 q_{\mathsf M(a,\cdot)}(z)
 =4r(a)r(z)+8\operatorname{Re}(D(a)\overline{D(z)}).
\]
On a common Hesse cubic, therefore, the vectors
\(v(z)=(r(z),\sqrt2\operatorname{Re}D(z),\sqrt2\operatorname{Im}D(z))\)
have pairwise nonnegative inner products. A second design has mean
\(v=0\). The nonnegative continuous kernel \(v(a)\cdot v(z)\) then has
zero double integral, forcing \(v=0\) on the entire support.

Suppose all coordinates of such a point are nonzero. For a unit
representative, put \(p_i=|z_i|^2\), \(u_i=z_i^3/p_i\).
The equations \(r=D=0\) imply
\[
 \sum_i u_i=0,\qquad |u_i|^2=p_i,\qquad \sum_i p_i^2=\frac12.
\]
Heron's identity gives zero area for the triangle of vectors \(u_i\).
Hence \(u_i=e^{i\theta}s_i\), where \(s_i\) are real,
\(\sum_i s_i=0\), and \(s_i^2=p_i\). Thus
\(z_i=e^{i\theta/3}\zeta_i s_i\), \(\zeta_i^3=1\), and
\[
 \frac{S(z)}{T(z)}=\frac3{\zeta_0\zeta_1\zeta_2}.
\]
This is a singular Hesse parameter. On a smooth member one coordinate
must therefore vanish; \(r=0\) equalizes the other two moduli, and the
cubic equation makes their cubes opposite. The support is contained
in \(\mathcal H\). Pairing \(M_2=I_6/6\) with \(P_i^{\otimes2}\)
gives
\[
 \frac16=w_i+\frac1{16}(1-w_i),
\]
so \(w_i=1/9\). Conversely those weights form a second design, by direct
summation or equality in \(\operatorname{tr}M_2^2\ge1/6\).

\emph{Step 5: the triangle.}
On \(xyz=0\), \(\sum_i|z_i|^4\ge1/2\), whose second-design expectation
is \(1/2\). Thus the support lies on the three midpoint circles.
The mixed moment \(\int|z_j|^2|z_k|^2\,d\mu=1/12\) gives mass \(1/3\)
to each circle. The off-diagonal identities are precisely vanishing
of its first and second phase Fourier coefficients. Conversely those
conditions give every bidegree-\((2,2)\) moment.
\end{proof}

\paragraph{Why this tensor format is special.}
For an \(n\times n\) pencil in \(d\) variables, the rank witness
\((n-1)\operatorname{tr}A^2-(\operatorname{tr}A)^2\) is nonnegative on
the determinant hypersurface. After balancing its first two factors,
a supported second design would imply
\[
 0\le\frac{n-2}{n}s^2-\operatorname{tr}T^2
 \le\left(\frac{n-2}{n}-\frac1d\right)s^2.
\]
The strictly negative range excludes designs. Equality forces the
third factor to balance, and for integers \(n,d\ge2\) occurs exactly at
\((n,d)=(3,3)\) and \((4,2)\). The binary-quartic case recovers a regular
tetrahedron on the Bloch sphere with uniform weights. Indeed, for a
second design on \(\mathbb P^1\), the nonnegative polynomial
\((3|\langle p,z\rangle|^2-1)^2\) has mean one and value four at \(p\),
so every atom has mass at most \(1/4\). A quartic has at most four
distinct zeros; equality therefore forces four equal weights, with
the polynomial zero at every other support point. Their squared
overlaps are \(1/3\), precisely the regular tetrahedron condition.
The three-qutrit case is a critical format of the rank inequality,
which explains the effectiveness of the exceptional-group description.

\subsection{An algebraic certificate that becomes a quantum experiment}
\label{geo:hesse-cert}

The cubic classification identifies an ensemble from exact equations.
The next result is quantitative and operational. A request of size \(k\)
means that a hidden ray \(z\sim\mu\) is drawn once and \(k\) copies of that
same state are prepared; independent requests use independent draws.
This is a matched-replica model with trusted dimension, basis, and
measurements. Its average request state is \(M_k(\mu)\). Redrawing the
hidden ray separately for every copy would instead reveal only \(M_1\).

Use pure-state trace distance
\[
 d(z,w)=\frac12\|P_z-P_w\|_1
       =\sqrt{1-|\langle z,w\rangle|^2}
\]
and its \(1\)-Wasserstein distance \(W_1\) on ensembles.
In \(\operatorname{Sym}^3\mathbb C^3\), define orthonormal vectors
\[
 \phi_F=\frac{|000\rangle+|111\rangle+|222\rangle}{\sqrt3},\qquad
 \phi_T=\frac1{\sqrt6}\sum_{\pi\in S_3}|\pi(012)\rangle,\qquad
 Q=|\phi_F\rangle\langle\phi_F|+|\phi_T\rangle\langle\phi_T|.
\]
Its coherent-state symbol is
\begin{equation}
 q(z)=\langle z^{\otimes3},Qz^{\otimes3}\rangle
 =\frac13|z_0^3+z_1^3+z_2^3|^2+6|z_0z_1z_2|^2.
 \label{geo:q}
\end{equation}
The two cubic equations cut out precisely the nine rays \(\mathcal H\).

\fig{fig02_cubic_witness}{A positive cubic witness locates the Hesse
ensemble. (a) The exact symbol $q$ on the slice
$z=(0,\cos s,e^{i\theta}\sin s)$ has three isolated zeros at equal
coordinate moduli and the Hesse phases. (b) For 2,500 independent complex
Gaussian rays normalized to unit length (fixed reproducible seed), the
distance to the nearest Hesse ray is plotted against $q$. The solid curve
is the proved bound $\min\{1,\sqrt{27q/4}\}$; the samples illustrate the
bound and do not establish it. The inset shows the same symbol near one
zero as a surface, with the plane $q=0.03$ for scale.}{fig:witness}

\begin{theorem}[Robust Hesse ensemble certificate]
\label{geo:hesse-bound}
For any Borel probability measure \(\mu\) on \(\mathbb P^2\), put
\[
 \epsilon_Q=\operatorname{tr}(QM_3(\mu)),\qquad
 \delta_1=\|M_1(\mu)-I_3/3\|_{\rm HS},\qquad K=1+\frac3{\sqrt2}.
\]
Then
\begin{equation}
 W_1(\mu,\mu_{\mathcal H})
 \le K\sqrt{\frac{27}{4}\epsilon_Q}+\frac32\delta_1.
 \label{geo:robust}
\end{equation}
In particular, \(M_1(\mu)=I_3/3\) and \(\epsilon_Q=0\) characterize
\(\mu_{\mathcal H}\) among all probability ensembles.
\end{theorem}

\begin{proof}
First we prove the global error bound
\begin{equation}
 d(z,\mathcal H)^2\le\frac{27}{4}q(z).
 \label{geo:pointwise}
\end{equation}
For a normalized binary vector \((u,v)\), let \(\Delta\le\pi/3\) be its
phase error from the nearest vector \((1,-\omega^j)/\sqrt2\).
Put \(t=2|uv|\), \(y=t\cos\Delta\). Its squared distance to this
three-point set is \((1-y)/2\), whereas
\[
 |u^3+v^3|^2
 =(1-y)\left(1+y+y^2-\frac34t^2\right).
\]
The second factor is at least one, since
\[
 y+y^2-\frac34t^2
 =(y-t/2)(y+t/2+1)+\frac12t(1-t)\ge0.
\]
Thus binary squared distance is at most \(|u^3+v^3|^2/2\).

For a qutrit let \(a\ge b\ge c\) be its squared coordinate moduli and
\(f=|\sum_i z_i^3|\). Delete the minimum coordinate and use the binary
estimate. Since \(c\le1/3\),
\[
 d(z,\mathcal H)^2
 \le c+\frac{(f+c^{3/2})^2}{2(1-c)^2}
 \le\frac54c+\frac94f^2.
\]
If \(a\le3/4\), then \(ab\ge3/32\), and this is at most
\(40abc/3+9f^2/4\le(27/4)q(z)\).
If \(a\ge3/4\), then
\[
 f\ge a^{3/2}-(1-a)^{3/2}\ge2a-1,\qquad
 d(z,\mathcal H)^2\le\frac{3-a}{4}.
\]
The identity
\(9(2a-1)^2-(3-a)=(4a-3)(9a-2)\ge0\) again proves
\eqref{geo:pointwise}.

Map each ray to a nearest point of \(\mathcal H\), obtaining weights
\(w_i\), and put \(L^2=\int d(z,\mathcal H)^2\,d\mu\).
This transport costs at most \(L\), and the first-moment error of the
transported measure is at most \(\delta_1+\sqrt2L\).
For \(u_i=w_i-1/9\), the SIC Gram matrix gives
\[
 \left\|\sum_i u_iP_i\right\|_{\rm HS}^2=\frac34\sum_i u_i^2.
\]
The weights' total-variation distance from uniform is at most
\(\sqrt3\,\delta_1+\sqrt6\,L\).
All pairwise Hesse distances are \(\sqrt3/2\), so correcting the weights
costs at most \((3/2)\delta_1+(3/\sqrt2)L\).
Combine the two transports and \(L^2\le27\epsilon_Q/4\).
\end{proof}

All measurements in the theorem are positive operators.
The Hesse third moment has the particularly simple form
\begin{equation}
 M_3(\mu_{\mathcal H})=\frac18(I_{10}-Q).
 \label{geo:hesse-third}
\end{equation}
The vectors \(h^{\otimes3}\) form a nine-vector equiangular tight frame
in the eight-dimensional kernel of \(Q\). Algebraically, the ideal
\((z_0^3+z_1^3+z_2^3,z_0z_1z_2)\) is a complete intersection, with
Hilbert series
\[
 \frac{(1-t^3)^2}{(1-t)^3}.
\]
Its Hilbert function starts \(1,3,6,8,9,9,\ldots\), giving the ranks
available to coherent tensors at successive replica numbers.
At three replicas the density matrix has rank eight, yet its
decomposition into pure coherent tensors is uniquely the nine-point
Hesse mixture. Indeed, the coherent mixtures annihilated by \(Q\) are
supported on \(\mathcal H\), and the Gram matrix of their nine
projectors is
\[
 \frac{63I_9+\mathbf1\mathbf1^{\mathsf T}}{64},
\]
so those projectors are linearly independent. They span an exposed
simplex face of the convex set of bosonic separable states.

Three matched replicas are also necessary for distribution-free
identification. Haar measure and \(\mu_{\mathcal H}\) have equal first
and second moments, so any protocol with independent request blocks of
size at most two receives identical quantum states in the two cases,
including joint processing of all blocks. At order three,
Haar has \(M_3=I_{10}/10\) and \(Q\)-probability \(1/5\), whereas Hesse
has \(Q\)-probability zero. The degree of the defining equations has
therefore acquired a direct experimental meaning.

\paragraph{Finite-sample certification.}
On \(N_3\) triple requests measure \(\{Q,I-Q\}\), observing \(s\) clicks.
Let \(q_U\) be the one-sided Clopper--Pearson upper bound at failure
probability \(\alpha/2\). For \(s<N_3\), it is the
\(1-\alpha/2\) quantile of \(\operatorname{Beta}(s+1,N_3-s)\); if
\(s=N_3\), put \(q_U=1\). In particular, zero clicks give
\[
 q_U=1-(\alpha/2)^{1/N_3}.
\]
On \(N_1\) single requests measure the SIC POVM \(\{P_h/3\}\),
with empirical frequencies \(\widehat p_h\), and set
\[
 \widehat\delta_1=\sqrt{12}
 \left\|\widehat p-\frac19\mathbf1\right\|_2,\qquad
 u_{N_1}=
 \frac{\sqrt{12}\bigl(1+\sqrt{\log(2/\alpha)}\bigr)}{\sqrt{N_1}}.
\]
Multinomial mean-square error and bounded differences imply, with
confidence at least \(1-\alpha\),
\begin{equation}
 W_1(\mu,\mu_{\mathcal H})
 \le\min\left\{1,\,
 K\sqrt{\frac{27}{4}q_U}
 +\frac32(\widehat\delta_1+u_{N_1})\right\}.
 \label{geo:confidence}
\end{equation}
The inverse SIC relation is \(M_1=4\sum_h p_hP_h-I\).
For example, the sufficient budgets
\[
 N_3\ge27K^2\varepsilon^{-2}\log(2/\alpha),\qquad
 N_1\ge432\varepsilon^{-2}
                 \bigl(1+\sqrt{\log(2/\alpha)}\bigr)^2
\]
make the acceptance rule \(s=0\),
\(\widehat\delta_1\le\varepsilon/6\) sound and complete for the ideal
Hesse ensemble at confidence \(1-\alpha\). The number of physical
copies is \(3N_3+N_1\).

The accuracy exponent \(\varepsilon^{-2}\) is optimal for fixed request
patterns. For \(0\le\theta\le\pi/6\), rotate the ensemble by
\(U_\theta=\operatorname{diag}(e^{i\theta},e^{-i\theta},1)\). Then
\[
 W_1(\mu_{\mathcal H},U_{\theta*}\mu_{\mathcal H})
 =\frac{\sin\theta+2\sin(\theta/2)}3,\qquad
 \epsilon_Q=
 \frac{\sin^2(3\theta)+2\sin^2(3\theta/2)}{18}.
\]
Coupling hidden labels and using tensor-product fidelity bounds the
trace distance of the two \(n\)-copy experiments by \(\sqrt n\sin\theta\).
Constant-success testing at ensemble distance \(\varepsilon\) thus
requires \(n=\Omega(\varepsilon^{-2})\). This assertion concerns
fixed request patterns; it makes no claim about every adaptive
stopping convention.

\paragraph{A portable principle.}
Let a finite ray set \(S=\{h_1,\ldots,h_N\}\) be cut out by degree-\(k\)
equations whose coefficient space has orthogonal projector \(Q\).
Suppose a certified global bound is available:
\[
 d(z,S)^2\le
 C\langle z^{\otimes k},Qz^{\otimes k}\rangle.
\]
Suppose also that the weight map
\(u\mapsto\sum_j u_jP_{h_j}^{\otimes r}\), restricted to
\(\sum_j u_j=0\), has least singular value \(s_r>0\). If the observed
\(r\)-th moment differs from that of target weights by \(\delta\)
in Hilbert--Schmidt norm, and the \(Q\)-probability is \(\epsilon\),
the same transport argument gives
\begin{equation}
 W_1(\mu,\mu_*)
 \le\sqrt{C\epsilon}
 +\frac{\operatorname{diam}(S)\sqrt N}{2s_r}
       \left(\delta+\sqrt{2rC\epsilon}\right).
 \label{geo:general-cert}
\end{equation}
Equations localize the support; moments condition its weights.
Transverse isolated zeros give local quadratic error bounds. An
effective global \(C\) additionally needs a Jacobian estimate near the
zeros and a positive minimum away from them. These are concrete
targets for exact algebra and certified numerics.

\subsection{A second design that obstructs geometric localization}
\label{geo:singer}

The first-moment proof used a coupling on well-conditioned projective
frames. A natural stronger localization at second moment would couple
nine points whose blowup has nef anticanonical class. Here \emph{nef}
means nonnegative intersection with every curve. The next example
shows that second-moment isotropy alone cannot enforce that geometry.

Let \(\zeta=e^{2\pi i/7}\), and set
\begin{equation}
 B=\{e_0,e_1,e_2\},\qquad
 T_j=\frac{(1,\zeta^j,\zeta^{3j})}{\sqrt3},\quad 0\le j<7,\qquad
 \mu_{\rm S}=\frac1{12}\sum_{b\in B}\delta_b+
                  \frac3{28}\sum_{j=0}^6\delta_{T_j}.
 \label{geo:singer-measure}
\end{equation}
This ten-point design is a known projective-toric construction
\cite[Example 4.2]{geo_iosue2024}. Its character exponents
\(\{0,2,6,1,3,4\}\) in the symmetric-square basis are distinct modulo
seven, giving \(M_2(\mu_{\rm S})=I_6/6\) by exact character cancellation.

\begin{theorem}[An exact coupling obstruction and repair threshold]
\label{geo:singer-nef}
There is no coupling with all marginals \(\mu_{\rm S}\) supported on
nine distinct points whose blowup of \(\mathbb P^2\) has nef
anticanonical class. For
\[
 \mu_t=t\mu_{\rm S}+(1-t)\sigma,
\]
such a coupling exists if and only if \(t\le20/21\).
The coupling may always be taken exchangeable.
\end{theorem}

\begin{proof}
The cubics
\[
 F_0=x^3-yz^2,\qquad F_1=y^3-x^2z,\qquad F_2=z^3-xy^2
\]
vanish at all \(T_j\). Each omits one basis point and has a cusp at
another. A nine-point subset obtained by deleting a basis point thus
has a cubic through eight points simply and one doubly; its strict
transform has anticanonical intersection \(9-10=-1\).
Deleting a torus point instead leaves nine points on a smooth cubic.
For deletion of \(T_0\), an example is
\[
 x^2y+x^2z+xy^2+y^2z+xz^2+yz^2+xyz=0.
\]
Its gradient differences are
\[
 G_x-G_y=(y-x)(x+y-z),\quad
 G_y-G_z=(z-y)(y+z-x),\quad
 G_z-G_x=(x-z)(z+x-y).
\]
At a common gradient zero, if \(x=y=z\), a gradient entry is \(7x^2\).
If exactly two coordinates agree, say \(x=y\ne z\), the remaining
differences force \(z=0\), but then \(G_z=3x^2\ne0\).
If all three are distinct, the differences force
\(x+y=z,\ y+z=x,\ z+x=y\), hence the zero vector.
The cubic is therefore smooth, and
the diagonal order-seven symmetry handles all other deletions.
A smooth cubic through nine points has a strict transform of square
zero; its intersections with every distinct irreducible curve are
nonnegative, so it is nef.

Thus a nef nine-subset of the ten-point support contains all three
basis points. Its expected number of basis points would be three,
whereas its prescribed marginals give \(9\mu_{\rm S}(B)=9/4\).

For a tuple with marginal \(\mu_t\), let \(n_T\) count entries in the
seven-point torus orbit, and \(n_O\) entries outside the ten-point
support. Every admissible tuple satisfies \(n_T\le6+n_O\): when
\(n_O=0\) it must omit a torus point, and otherwise distinctness gives
the assertion. Taking expectations yields
\[
 \frac{27t}{4}\le6+9(1-t),\qquad t\le\frac{20}{21}.
\]
At equality, mix two types and randomly permute the tuple:
with probability \(4/7\), all three basis points and six uniformly
chosen torus points; with probability \(3/7\), all seven torus points,
one uniform basis point, and one Haar point.
For the second type, cyclic symmetry reduces to the retained basis
point \(e_0\). Cubics through \(T\cup\{e_0\}\) form the pencil
\(\operatorname{span}(F_1,F_2)\), since the seven cubic evaluation
characters exhaust all residues modulo seven.
This pencil contains the smooth member
\[
 F_1+F_2=y^3-x^2z+z^3-xy^2.
\]
Indeed, its \(y\)-derivative is \(y(3y-2x)\). If \(y=0\), the other
gradient equations force the zero vector. Otherwise \(y=2x/3\),
the \(x\)-derivative forces \(z=-2x/9\), and the \(z\)-derivative is
\((4/27-1)x^2\ne0\). Only finitely many pencil members are singular,
so a Haar point chooses a smooth member almost surely.
The mixture has weights \(5/63\) at each basis point, \(5/49\) at
each torus point, and Haar weight \(1/21\), exactly \(\mu_{20/21}\).
For smaller \(t\), mix with independent Haar points, whose generic
nine-point configuration lies on a smooth cubic.
\end{proof}

Two measures with identical second moments can therefore have
different coupling geometry. The obstruction is not explained by a
large atom: every atom of \(\mu_{\rm S}\) has mass below \(1/9\).
It is encoded in the incidence and multiplicities of the three cubics.

\subsection{The same design produces a singular K3 surface}
\label{geo:k3}

Blow up the ten points and call the rational surface \(S\). The strict
transforms \(C_0,C_1,C_2\) of the preceding cubics are disjoint smooth
rational curves with
\[
 C_i^2=-3,\qquad C_iC_j=0\ (i\ne j),\qquad
 C_0+C_1+C_2=-3K_S,\qquad K_S^2=-1.
\]
The cyclic triple cover \(\pi:Z\to S\), using the line bundle
\(-K_S\) and branch divisor \(\sum C_i\), is smooth: locally it has
equation \(w^3=u\). Its ramification curves \(R_i\) have square
\(-1\), and the canonical bundle formula gives
\(K_Z=R_0+R_1+R_2\). Contract them to obtain \(X\).
Then \(K_X\) is trivial, and
\[
 e(Z)=3e(S)-2\sum_ie(C_i)=3\cdot13-2\cdot6=27,\qquad e(X)=24.
\]
Noether's formula gives \(\chi(\mathcal O_X)=2\); connectedness and
triviality of \(K_X\) imply \(h^1(\mathcal O_X)=0\).
Thus \(X\) is a smooth projective K3 surface.
The term \emph{singular K3} below means maximal Picard rank twenty.

\fig{fig05_singer_cover}{From a finite projective design to a K3 surface.
(a) The three coordinate rays and seven Singer rays are drawn as an
incidence schematic; their positions in the drawing are not a projective
embedding. (b) Blowing up the ten rays separates the three cubic transforms
into disjoint $(-3)$-curves. The triple cover makes their ramification
curves $(-1)$-curves; contracting these gives a smooth K3 surface with
Euler number 24. (c) The transcendental lattice has the Eisenstein shape
$A_2(7)$, Gram matrix $\left(\begin{smallmatrix}14&-7\\-7&14\end{smallmatrix}\right)$,
and discriminant 147. The arrows in (b) describe actual geometric
operations, while (c) is a lattice diagram.}{fig:k3}

\begin{theorem}[The lattice of the design cover]
\label{geo:k3-lattice}
The K3 surface obtained from this ten-point configuration satisfies
\[
 \rho(X)=20,\qquad
 T_X\cong A_2(7)=
 \begin{pmatrix}14&-7\\-7&14\end{pmatrix},\qquad
 \operatorname{disc}T_X=147.
\]
It admits a symplectic automorphism group \(C_7\rtimes C_3\) and a
commuting non-symplectic deck transformation of order three with
exactly three isolated fixed points.
\end{theorem}

\begin{proof}
A weighted-projective model is
\[
 w^3=(x^3-yz^2)(y^3-x^2z)(z^3-xy^2),\qquad
 \deg(x,y,z,w)=(1,1,1,3).
\]
On its smooth locus a two-form is
\[
 \Omega=\frac{x\,dy\wedge dz-y\,dx\wedge dz+z\,dx\wedge dy}{w}.
\]
At each blown-up point the total branch multiplicity is three, so
the blowup removes its apparent poles. The zeros are exactly the
ramification curves, removed by their contraction.
The deck map \(\delta:w\mapsto e^{2\pi i/3}w\) acts nontrivially on
\(\Omega\); its three fixed curves become three isolated fixed points.
The maps
\[
 g:(x,y,z,w)\mapsto(x,\zeta y,\zeta^3z,\zeta^4w),\qquad
 c:(x,y,z,w)\mapsto(y,z,x,w)
\]
preserve \(\Omega\), and satisfy \(g^7=c^3=1\),
\(cgc^{-1}=g^2\) in weighted projective coordinates.

A symplectic order-seven automorphism has rank-four invariant
cohomology and a rank-eighteen algebraic coinvariant lattice
\cite[Propositions 4.5--4.6]{geo_garbagnati_sarti2007}.
An invariant ample class therefore gives \(\rho(X)\ge19\).
The non-symplectic order-three action makes the transcendental
lattice a module over \(\mathbb Z[e^{2\pi i/3}]\), so its rank is
even. It is at most three and at least two, hence is two.
Thus \(\rho(X)=20\). Since the Eisenstein ring is principal, its
invariant positive even rank-two form is \(A_2(m)\) for an integer \(m>0\).

To determine \(m\), let
\[
 L_0=\{D\in\operatorname{Pic}(S):D C_i=0\text{ for every }i\}.
\]
The primitive closure of \(\langle C_0,C_1,C_2\rangle\) has index
three: it contains \(K_S=-(C_0+C_1+C_2)/3\), and its index squared
divides \(27\). Unimodularity of \(\operatorname{Pic}(S)\) gives
\(|\operatorname{disc}L_0|=3\).
Pullback scales its form by three. The fixed lattice of \(\delta\)
has rank eight and discriminant \(3^7\), since its fixed locus is
three isolated points \cite[Theorem 4.1]{geo_artebani_sarti_taki2011}.
Hence \(\pi^*L_0\) has index three in that fixed lattice.

Let \(H\) and \(E_b,E_{T_j}\) be the line and exceptional classes on \(S\).
The \(g\)-fixed part of \(L_0\) has basis
\[
 v_1=H-\sum_{b\in B}E_b,\qquad
 v_2=7H-3\sum_{j=0}^6E_{T_j},\qquad
 \bigl(v_iv_j\bigr)=\begin{pmatrix}-2&7\\7&-14\end{pmatrix}.
\]
The order-seven action on the index-three quotient is trivial, and
the obstruction in \(H^1(C_7,\pi^*L_0)\) is seven-torsion. Thus
taking \(g\)-invariants preserves index three.
The unique even index-three overlattice of the threefold-scaled
matrix adds \((v_2-v_1)/3\): an order-three discriminant element
\((\alpha v_1+\beta v_2)/3\) is isotropic exactly when
\(\alpha+\beta=0\) modulo three.
In the basis \(v_1+(v_2-v_1)/3,-(v_2-v_1)/3\) its matrix is
\[
 N=\begin{pmatrix}2&1\\1&-10\end{pmatrix}.
\]
A symplectic automorphism fixes \(T_X\) pointwise, since its rational
coinvariant classes are of type \((1,1)\).
The rank-four invariant lattice therefore has rank-two algebraic part.
This contains the rank-two lattice just computed in the \(\delta\)-fixed
part; consequently \(\delta\) acts trivially on its rational span, and
\(\operatorname{NS}(X)^g=N\).

The full invariant lattice is
\[
 H^2(X,\mathbb Z)^g
 \cong U(7)\oplus\begin{pmatrix}4&1\\1&2\end{pmatrix},
 \qquad |\operatorname{disc}|=7^3,
\]
with elementary discriminant group \((\mathbb Z/7)^3\)
\cite[Proposition 4.6]{geo_garbagnati_sarti2007}.
Both \(N\) and \(T_X\) are primitive in this lattice.
If \(h\) is their gluing index, the gluing group injects into both
discriminant groups, so \(h\mid21\). Determinants give
\[
 7^3=\frac{21\cdot3m^2}{h^2}.
\]
The only possibilities are \((m,h)=(7,3),(49,21)\).
In the latter case, the seven-primary discriminants of \(N,T_X\)
are \(\mathbb Z/7\) and \((\mathbb Z/49)^2\). The gluing subgroup
has generator \((n,7v)\) with \(n\ne0\).
Its orthogonal complement projects onto \((\mathbb Z/49)^2\),
since every pairing with \(7v\) can be canceled by a multiple of \(n\).
Lift an element of order \(49\). Seven times that lift has zero first
component and nonzero second component; it cannot belong to the
gluing subgroup. The glued discriminant would retain an element of
order \(49\), contradicting its elementary form. Therefore \(m=7\).
\end{proof}

The invariant lattice also makes the relevant nef geometry explicit.
Its norm is \(2(x^2+xy-5y^2)\), congruent to \(2(x-y)^2\) modulo three,
so it represents neither \(-2\) nor \(-14\).
A rational curve on \(X\) has square \(-2\); its \(g\)-orbit cannot
have length one. A length-seven orbit sum has square \(-14+14k\)
for an integer \(k\ge0\), and cannot have square \(-14\), hence is
nonnegative. Classes in the positive component of \(N_{\mathbb R}\)
therefore pair positively with every such orbit sum, and thus with
every curve, by the Hodge index theorem. The invariant nef cone is
the closure of that positive component. In invariant blowup
coordinates
\[
 dH-b\sum_{b\in B}E_b-c\sum_jE_{T_j},\qquad
 3d=3b+7c,
\]
its null boundary rays are \(d/b=(9\pm\sqrt{21})/2\).

The construction does not assert a new K3 isomorphism class. Its
content is the explicit passage from an isotropic quantum ensemble,
through a geometric localization obstruction, to a canonical cover
whose transcendental lattice is determined. The same cubic
multiplicities that obstruct nef localization make the cyclic cover
possible.

\subsection{Verification and reproducible mathematical content}
\label{geo:verification}

The conceptual inputs are Kempf--Ness, the three-qutrit Cartan
classification, semicontinuity of singularity exponents,
Hilbert/Bergman geodesic comparison, and the stated K3 automorphism
lattices. They enter with precise hypotheses and primary references.
Finite calculations are supplied as exact-arithmetic certificates:
Moore determinants and quartic contractions; Hesse moments, Gram
matrices and spectra; Singer character cancellation, cubic
smoothness and intersections; and finite minimal-jet moment-map
examples. Numerical sampling illustrates bounds and searches for
counterexamples. The theorem constants follow from the stated
inequalities and identities.

The operational supplement evaluates \eqref{geo:confidence} from
measurement counts and supplies seeded synthetic-data examples.
It exposes the equation-to-projector and moment-to-weight-conditioning
steps of \eqref{geo:general-cert} for reuse with other finite ray
sets. The computations verify finite identities and statistical
calculations; they are not a completed machine formalization of
the analytic and geometric proofs.

